\SGAchapter{III}{光滑态射：延拓性质}{Morphismes lisses : propriétés de prolongement}
\SourcePage{65}
\SGAsection{1}{形式光滑同态}
在第 II 讲中，我们只考虑了\Term{finite-type}{有限型}\Term{homomorphism}{同态}；因此，对于\Term{local-homomorphism}{局部环同态} $A\to B$，也只考虑了 $B$ \Term{isomorphism}{同构}于某个\Term{finite-type}{有限型} $A$-代数的\Term{localization}{局部化}的情形。这不足以满足各种应用的需要，特别是在\Term{formal-geometry}{形式几何}或解析几何中。例如，形式幂级数环 $B=A[[t_1,\ldots,t_n]]$，从\Term{formal-geometry}{形式几何}的观点看，具有 $A$ 上\Term{smooth}{光滑}代数的性质。在解析几何中，乘积 $Y\times\mathbf C^n$ 的一点 $(y,z)$ 的\Term{local-ring}{局部环}，视为 $y$ 的\Term{local-ring}{局部环}上的代数时，也具有这一性质；而且，这个代数的\Term{completion}{完备化}\Term{isomorphism}{同构}于基环 $O_x$ 的\Term{completion}{完备化}上的 $n$ 元\Term{formal-power-series-algebra}{形式幂级数代数}。这就引出了以下定义。
\begin{SGAresult}{定义}{1.1}
设 $u:A\to B$ 是\Term{local-ring}{局部环}之间的\Term{local-homomorphism}{局部同态}（提醒读者：环均为\Term{noetherian-ring}{诺特环}），且 $k(B)$ 在 $k(A)$ 上\Term{finite}{有限}。若存在一个\Term{finite}{有限}的局部 $\widehat A$-代数 $A'$，它在 $\widehat A$ 上自由，并使半\Term{local-ring}{局部环} $B'=\widehat B\otimes_{\widehat A}A'$ 的各局部分量均 $A'$-\Term{isomorphism}{同构}于 $A'$ 上的\Term{formal-power-series-algebra}{形式幂级数代数}，则称 $u$ 为\Term{formally-smooth}{形式光滑同态}，或称代数 $B$ 在 $A$ 上\Term{formally-smooth}{形式光滑}。\footnote{关于一个由下文 2.1 引出的、更一般而且更具概念性的定义，参见 EGA $0_{\mathrm{IV}}$ 19.3.1。}
\end{SGAresult}
（$\widehat A,\widehat B$ 表示 $A,B$ 的\Term{completion}{完备化}。）由于 $B'$ 在 $\widehat B$ 上\Term{finite}{有限}且自由，它确实是半\Term{local-ring}{局部环}，是若干完备\Term{local-ring}{局部环}的直积；每个分量仍为自由 $\widehat B$-模，因而维数等于 $\widehat B$，也等于 $B$ 的维数。所以，1.1 中各形式幂级数环的变量个数等于
\[\dim\widehat B-\dim\widehat A=\dim B-\dim A,\]
特别地，它与所选的局部分量无关。立刻可见，这也是环 $B\otimes_A k=B/\mathfrak mB$ 的维数，其中 $k=A/\mathfrak m$ 为 $A$ 的\Term{residue-field}{剩余域}；称之为 $B$ 相对于 $A$ 的\Term{relative-dimension}{相对维数}。
\SourcePage{66}
\textbf{注记 1.2。}显然，定义 1.1 只依赖于 $A\to B$ 所诱导的\Term{completion}{完备化}\Term{homomorphism}{同态} $\widehat A\to\widehat B$，这在一定程度上说明了这一名称的合理性。这里我们要对第 I 讲 3.2 b) 和 4.1 b) 的定义表示遗憾：那些定义容易引起误解。在那些情形中，我们宁愿说“\Term{formally-unramified}{形式非分歧}”和“\Term{formally-etale}{形式平展}”，而把“\Term{unramified}{非分歧}”和“\Term{etale}{平展}”留给 $B$ 是\Term{finite-type}{有限型} $A$-代数的\Term{localization}{局部化}的情形。\footnote{更好的是说“本质\Term{unramified}{非分歧}”和“本质\Term{etale}{平展}”，比较 EGA IV 18.6.1。}读者立刻可以验证：“\Term{formally-etale}{形式平展}”等价于“\Term{formally-smooth}{形式光滑}且\Term{quasi-finite}{拟有限}”。最后指出，不对\Term{residue-field}{剩余域}\Term{group-extension}{扩张} $k(B)/k(A)$ 预加任何条件（这里假定它\Term{finite}{有限}），也能合理定义“\Term{formally-smooth}{形式光滑}”；这个定义尤其包含如下\Term{local-homomorphism}{局部同态} $A\to B$：$B$ 在 $A$ 上\Term{flat}{平坦}，而且 $B/\mathfrak mB$ 是 $A/\mathfrak m=k$ 的\Term{separable-extension}{可分扩张}（不一定\Term{finite-type}{有限型}）。例如，Cohen 的 $p$-环在 $p$ 进整数环上\Term{formally-smooth}{形式光滑}。在这一一般情形中，应当把\Term{homomorphism}{同态}的提升性质（比较 2.1）作为定义。对于我们所考虑的应用，定义 1.1 的情形已经足够；以下说“\Term{formally-smooth}{形式光滑}”时，均默认为“\Term{residue-field}{剩余域}\Term{group-extension}{扩张}\Term{finite}{有限}”。
\begin{SGAresult}{引理}{1.3}
若 $B$ 在 $A$ 上\Term{formally-smooth}{形式光滑}，则 $B$ 在 $A$ 上\Term{flat}{平坦}。
\end{SGAresult}
\begin{proof}
由于\Term{flat}{平坦}性在\Term{completion}{完备化}下不变，可以假定 $A,B$ 完备。由于\Term{flat}{平坦}性在基环的局部\Term{flat}{平坦}\Term{group-extension}{扩张}（因而是忠实\Term{flat}{平坦}\Term{group-extension}{扩张}）下不变，根据定义 1.1，只需考虑 $B$ 为 $A$ 上\Term{formal-power-series-algebra}{形式幂级数代数}的情形。但此时，作为 $A$-模，$B$ \Term{isomorphism}{同构}于一些\Term{isomorphism}{同构}于 $A$ 的 $A$-模的直积，因此（基环 $A$ 为\Term{noetherian-ring}{诺特环}）它作为\Term{flat}{平坦} $A$-模的直积也是 $A$-\Term{flat}{平坦}的。
\end{proof}
在 1.1 的条件下，$B'$ 的各局部分量在 $A'$ 上的\Term{residue-field}{剩余域}\Term{group-extension}{扩张}均\Term{trivial-bundle}{平凡}，故 $L\otimes_k k'$ 是一个 Artin $k'$-代数，其各局部分量的\Term{residue-field}{剩余域}\Term{group-extension}{扩张}均\Term{trivial-bundle}{平凡}（$L,k,k'$ 是 $A,B,A'$ 的\Term{residue-field}{剩余域}）。对于\Term{finite}{有限}自由\Term{group-extension}{扩张} $A'$ 满足 1.1 中的条件，这是必要条件；由下文 1.4 (i) 和 1.5 立即可见，它也充分。
\begin{SGAresult}{命题}{1.4}
设 $A\to B$ 是\Term{local-ring}{局部环}之间的\Term{local-homomorphism}{局部同态}，\Term{residue-field}{剩余域}\Term{group-extension}{扩张}\Term{finite}{有限}；$A'$ 为\Term{finite}{有限}的局部 $A$-代数，于是 $B'=B\otimes_A A'$ 在 $B$ 上\Term{finite}{有限}，因而也是\Term{noetherian}{诺特}半\Term{local-ring}{局部环}。(i) 若 $B$ 在 $A$ 上\Term{formally-smooth}{形式光滑}，则 $B$ 在其\Term{maximal-ideal}{极大理想}处的\Term{localization}{局部化}在 $A'$ 上\Term{formally-smooth}{形式光滑}。(ii) 若 $A'$ 在 $A$ 上自由，则逆命题成立。
\end{SGAresult}
\SourcePage{67}
\begin{proof}
立即归结到 $A,B$ 完备的情形。
(i) 取 $A$ 的\Term{finite}{有限}自由局部\Term{group-extension}{扩张} $A''$，使 $B''=B\otimes_A A''$ 的各局部分量都是 $A''$ 上的\Term{formal-power-series-algebra}{形式幂级数代数}。作标量\Term{group-extension}{扩张} $A''\to A''\otimes_A A'\to A'''$，其中 $A'''$ 为 $A''\otimes_A A'$ 的一个局部分量，便知 $B''\otimes_{A''}A'''=B\otimes_A A'''$ 的各局部分量均为 $A'''$ 上的\Term{formal-power-series-algebra}{形式幂级数代数}。另一方面，
\[B\otimes_A A'''=(B\otimes_A A')\otimes_{A'}A'''=B'\otimes_{A'}A'''.\]
又因 $A''$ 在 $A$ 上自由，$A''\otimes_A A'$ 在 $A'$ 上自由，作为其\Term{direct-summand}{直和因子}的 $A'''$ 也如此。这就证明 $B'$ 在 $A'$ 上\Term{formally-smooth}{形式光滑}。
(ii) 取\Term{finite}{有限}自由局部 $A'$-代数 $A''$，使 $B'\otimes_{A'}A''=B\otimes_A A''$ 的各局部分量为 $A''$ 上的\Term{formal-power-series-algebra}{形式幂级数代数}。由于 $A'$ 在 $A$ 上自由，$A''$ 亦然，所以 $B$ 在 $A$ 上\Term{formally-smooth}{形式光滑}。
\end{proof}
\begin{SGAresult}{命题}{1.5}
设 $A\to B$ 是\Term{local-ring}{局部环}之间的\Term{local-homomorphism}{局部同态}，\Term{residue-field}{剩余域}\Term{group-extension}{扩张}\Term{trivial-bundle}{平凡}。则 $B$ 在 $A$ 上\Term{formally-smooth}{形式光滑}的充要条件是：$\widehat B$ \Term{isomorphism}{同构}于 $\widehat A$ 上的\Term{formal-power-series-algebra}{形式幂级数代数}。
\end{SGAresult}
\begin{proof}
只需证明必要性，并可假定 $A,B$ 完备。设 $\mathfrak m$（$\mathfrak n$）为 $A$（$B$）的\Term{maximal-ideal}{极大理想}，取 $\mathfrak n$ 中的元素 $t_1,\ldots,t_n$，其像构成向量空间
\[(\mathfrak n/\mathfrak n^2)/\operatorname{Im}(\mathfrak m/\mathfrak m^2)=\mathfrak n/(\mathfrak n^2+\mathfrak mB)\]
的一组基。它们定义局部 $A$-代数\Term{homomorphism}{同态} $B_1=A[[t_1,\ldots,t_n]]\to B$；证明这是\Term{isomorphism}{同构}。只需证明对 $\mathfrak m$ 的任意幂 $\mathfrak m^q$，模 $\mathfrak m^q$ \Term{reduced}{约化}后得到\Term{isomorphism}{同构}，因为 $B_1,B$ 是相应\Term{reduced}{约化}环随 $q$ 变化所得的\Term{projective-limit}{逆极限}。由于 $B,B_1$ 是\Term{flat}{平坦} $A$-模，其 $\mathfrak m$ 进滤过的伴随\Term{grading}{分次}环，分别由 $B/\mathfrak mB,B_1/\mathfrak mB_1$ 在 $k=A/\mathfrak m$ 上与 $\operatorname{gr}(A)$ 张量得到。因此只需证明 $B_1/\mathfrak mB_1\to B/\mathfrak mB$ 是\Term{isomorphism}{同构}。结合 1.3，便归结到 $A$ 为域 $k$ 的情形。另一方面，取\Term{finite}{有限}自由局部 $A$-代数 $A'$，使 $B\otimes_A A'$ 为 $A'$ 上\Term{formal-power-series-algebra}{形式幂级数代数}（注意：因 $B/A$ 的\Term{residue-field}{剩余域}\Term{group-extension}{扩张}\Term{trivial-bundle}{平凡}，这个代数是局部的）。要证明 $B_1\to B$ 是\Term{isomorphism}{同构}，只需证明 $B_1\otimes_A A'\to B\otimes_A A'$ 是\Term{isomorphism}{同构}。因此归结到 $B$ 已是\Term{formal-power-series-algebra}{形式幂级数代数}的情形（应当在化到基环是域以前先作这一步）。此时 $B$ 是具有代表域 $k$ 的\Term{regular}{正则}\Term{local-ring}{局部环}，而由 $B_1,B$ 的 $\mathfrak n_1$ 进、$\mathfrak n$ 进滤过的伴随\Term{grading}{分次}环立即可见，$B_1\to B$ 为\Term{isomorphism}{同构}。这也是熟知的事实，证明完毕。
\end{proof}
\SourcePage{68}
\begin{SGAresult}{推论}{1.6}
若 $B$ 在 $A$ 上\Term{formally-smooth}{形式光滑}，则存在\Term{finite}{有限}局部 $A$-代数 $A'$，使 $\widehat B\otimes_{\widehat A}\widehat {A'}=\widehat{(B\otimes_A A')}$ 的各局部分量\Term{isomorphism}{同构}于 $\widehat {A'}$ 上的\Term{formal-power-series-algebra}{形式幂级数代数}。
\end{SGAresult}
\begin{proof}
设 $L/k$ 为 $B/A$ 的\Term{residue-field}{剩余域}\Term{group-extension}{扩张}，取\Term{group-extension}{扩张} $k'/k$，使 $k'$-代数 $L\otimes_k k'$ 的各\Term{residue-field}{剩余域}\Term{group-extension}{扩张}\Term{trivial-bundle}{平凡}。取\Term{finite}{有限}自由 $A$-代数 $A'$，使 $A'/\mathfrak mA'=k'$（这样的代数确实存在，例如逐步归结到 $k'/k$ 为单\Term{group-extension}{扩张}，再把 $k'$ 在 $k$ 上一个生成元的最小多项式的系数提升到 $A$）。它是局部的。于是 $B\otimes_A A'$ 在各\Term{maximal-ideal}{极大理想}处相对于 $A'$ 的\Term{residue-field}{剩余域} $k'$ 的\Term{group-extension}{扩张}均\Term{trivial-bundle}{平凡}，再用 1.5 即可。
\end{proof}
\begin{SGAresult}{推论}{1.7}
设 $A\to B$ 为\Term{local-ring}{局部环}之间的\Term{local-homomorphism}{局部同态}。$B$ 在 $A$ 上\Term{formally-smooth}{形式光滑}的充要条件是：$B$ 在 $A$ 上\Term{flat}{平坦}，且 $B/\mathfrak mB$ 在 $k=A/\mathfrak m$ 上\Term{formally-smooth}{形式光滑}。
\end{SGAresult}
\begin{proof}
作 $A$ 的适当\Term{finite}{有限}自由局部\Term{group-extension}{扩张} $A'$，由 1.4 (ii) 可归结到 $B/A$ 的\Term{residue-field}{剩余域}\Term{group-extension}{扩张}\Term{trivial-bundle}{平凡}的情形。由 1.4 (i) 和 1.3，所述条件必要。至于充分性，只要注意 1.5 的证明在这里的假设下也说明 $B$ 是 $A$ 上\Term{formal-power-series-algebra}{形式幂级数代数}即可（可假定 $A,B$ 完备）。
\end{proof}
\textbf{注记 1.8。}不难为\Term{formally-smooth}{形式光滑}\Term{homomorphism}{同态}建立第 II 讲中所讨论的\Term{smooth}{光滑}\Term{morphism}{态射}的各项性质的对应版本。但讨论微分性质时，需修改通常的 Kähler 微分定义（参见第 I 讲 1），以完备\Term{tensor-product}{张量积}代替通常\Term{tensor-product}{张量积}。这里仅提及这些深层问题，以上内容已足以供我们使用。还需把\Term{formally-smooth}{形式光滑}性与第 II 讲建立的\Term{smooth}{光滑}性概念联系起来（至今尚未用到后者）：
\begin{SGAresult}{命题}{1.9}
设 $A\to B$ 为\Term{local-homomorphism}{局部同态}，$B$ 为\Term{finite-type}{有限型} $A$-代数的\Term{localization}{局部化}。则 $B$ 在 $A$ 上\Term{smooth}{光滑}当且仅当它在 $A$ 上\Term{formally-smooth}{形式光滑}。
\end{SGAresult}
\begin{proof}
由 1.7 和第 II 讲 2.1，归结到 $A$ 为域的情形。利用 1.4 (ii) 和第 II 讲 4.13，作 $k$ 的适当\Term{group-extension}{扩张} $k'$，便可假定 $B/k$ 的\Term{residue-field}{剩余域}\Term{group-extension}{扩张}\Term{trivial-bundle}{平凡}。由 1.5（分别由第 II 讲 5.2），$B$ 在 $k$ 上\Term{smooth}{光滑}（分别为\Term{formally-smooth}{形式光滑}）的充要条件是 $B$ 为\Term{regular}{正则}\Term{local-ring}{局部环}（分别为其\Term{completion}{完备化}是 $k$ 上\Term{formal-power-series-algebra}{形式幂级数代数}）。而在\Term{residue-field}{剩余域}\Term{group-extension}{扩张}\Term{trivial-bundle}{平凡}时，这两个条件等价是熟知的。
\end{proof}
\SourcePage{69}
\SGAsection{2}{刻画形式光滑同态的提升性质}
\begin{SGAresult}{定理}{2.1}
设 $A\to B$ 为\Term{local-ring}{局部环}之间的\Term{local-homomorphism}{局部同态}，\Term{residue-field}{剩余域}\Term{group-extension}{扩张}\Term{finite}{有限}。下列条件等价：
\begin{enumerate}
\item[(i)] $B$ 在 $A$ 上\Term{formally-smooth}{形式光滑}。
\item[(ii)] 对任意\Term{local-homomorphism}{局部同态} $A\to C$（$C$ 为完备\Term{local-ring}{局部环}）、包含于 $r(C)$ 的任意理想 $J$，及任意局部 $A$-\Term{homomorphism}{同态} $B\to C/J$，存在一个提升它的 $A$-\Term{homomorphism}{同态} $B\to C$（必为\Term{local-homomorphism}{局部同态}）。
\item[(iii)] 对任意 $A$-代数 $C$（不必为\Term{noetherian-ring}{诺特环}）、$C$ 的任意\Term{nilpotent}{幂零}理想 $J$，及任意连续 $A$-\Term{homomorphism}{同态} $B\to C$（即在 $r(B)$ 的某个幂上为零），存在提升它的 $A$-\Term{homomorphism}{同态} $B\to C$（也必连续）。
\item[(iv)] 与 (ii)、(iii) 相同，但限 $C$ 为在 $A$ 上\Term{finite}{有限}的 Artin \Term{local-ring}{局部环}。
\item[(iv bis)] 与 (iv) 相同，并另要求 $J^2=0$。
\end{enumerate}
\end{SGAresult}
\textbf{注记。}本讲下文只会用到 (iv)$\Rightarrow$(i) 或 (iv bis)$\Rightarrow$(i)；当 $B$ 为\Term{finite-type}{有限型} $A$-代数的\Term{localization}{局部化}时，下一节将用另一方法证明正向蕴涵 (i)$\Rightarrow$(ii)。回顾 Cohen 定理的“正确”理论\footnote{参见 EGA $0_{\mathrm{IV}}$ 19.3、19.8。}，性质 (ii) 或 (iii) 成为\Term{formally-smooth}{形式光滑}\Term{homomorphism}{同态}的定义，而 1.1 则成为只适用于\Term{finite}{有限}\Term{residue-field}{剩余域}\Term{group-extension}{扩张}的刻画性质。注意 (ii)、(iii) 互不比对方更一般。若引入线性拓扑化、\Term{separated}{分离}且完备的环 $C$、$C$ 的拓扑\Term{nilpotent}{幂零}闭理想，以及连续\Term{homomorphism}{同态} $A\to C$（使 $C$ 成为拓扑 $A$-代数），便可给出一个同时包含两者的等价性质。这一修改留给读者。
\begin{proof}[2.1 的证明]
我们证明 (i)$\Rightarrow$(iii)$\Rightarrow$(ii)，再证明 (iv)$\Rightarrow$(i)。因 (ii)$\Rightarrow$(iv) 显然，而对满足 $J^n=0$ 的整数 $n$ 作直接归纳即可证明 (iv) 与 (iv bis) 等价，故这就完成证明。
(i)$\Rightarrow$(iii)。直接归纳可归结到 $J^2=0$。由于 $C$ 在 $A$ 上\Term{finite}{有限}，$A$ 的\Term{maximal-ideal}{极大理想}的某个幂 $\mathfrak m^q$ 消去 $C$。模去 $\mathfrak m^q$，并由 1.4 (i) 注意到 $B/\mathfrak m^qB$ 仍在 $A/\mathfrak m^q$ 上\Term{formally-smooth}{形式光滑}，可假定 $A$ 为 Artin 环。由 1.3，$B$ 在 $A$ 上\Term{flat}{平坦}；因 $A$ 为 Artin 环，$B$ 在 $A$ 上自由。于是存在 $A$-模\Term{homomorphism}{同态} $w:B\to C$，提升给定\Term{homomorphism}{同态} $u:B\to C/J$。
\SourcePage{70}
置 $f(x,y)=w(xy)-w(x)w(y)$（$x,y\in B$）。于是 $f(x,y)\in J$，且 $f:B\times B\to J$ 为 $A$-双线性映射。存在作为代数\Term{homomorphism}{同态}的提升 $v:B\to C$，当且仅当存在 $A$-线性映射 $g:B\to J$，使 $v=w+g$ 为代数\Term{homomorphism}{同态}；写成方程就是
\[g(1)=1-w(1),\qquad g(x,y)-u(x)g(y)-u(y)g(x)=-f(x,y)\quad(x,y\in B).\]
这是 $\operatorname{Hom}_A(B,J)$ 中的线性方程组，常数项在 $J$ 中。因此它有解当且仅当对于忠实\Term{flat}{平坦} $A$-代数 $A'$，$\operatorname{Hom}_A(B,J)\otimes_A A'$ 中常数项在 $J'=J\otimes_A A'$ 中的相应方程组有解。取\Term{finite}{有限}自由局部 $A$-代数 $A'$，使 $B'=B\otimes_A A'$ 为 $A'$ 上\Term{formal-power-series-algebra}{形式幂级数代数}（注意：立即可见，在此证明中可假定 $A,B$ 完备）。因 $A'$ 为\Term{finite-type}{有限型}自由 $A$-模，
\[\operatorname{Hom}_A(B,J)\otimes_A A'=\operatorname{Hom}_{A'}(B',J').\]
而 $\operatorname{Hom}_{A'}(B',J')$ 中的方程组，恰好决定如下 $A'$-代数\Term{homomorphism}{同态} $B'\to C'=C\otimes_A A'$：它们通过以 $A'$-模\Term{homomorphism}{同态} $g':B'\to J'$ “修正”由 $w$ 作标量\Term{group-extension}{扩张}所得的 $A'$-模\Term{homomorphism}{同态} $w':B'\to C'$，从而提升由 $u$ 作标量\Term{group-extension}{扩张}所得的 $u':B'\to C'/J'$。（注意 $B$ 作为 $A'$-模生成 $B'$。）因此只需在 $B=A[[t_1,\ldots,t_n]]$ 时证明 (iii)。任意把各 $t_i$ 在 $C/J$ 中的像提升为 $z_i\in C$。由于 $z_i\bmod J$ \Term{nilpotent}{幂零}（$u:B\to C/J$ 连续），且 $J$ \Term{nilpotent}{幂零}，$z_i$ 也\Term{nilpotent}{幂零}。因此 $z_i$ 定义从 $B$ 到\Term{discrete}{离散}环 $C$ 的连续拓扑 $A$-代数\Term{homomorphism}{同态}，显然提升 $u$。
(iii)$\Rightarrow$(ii)。设 $\mathfrak n$ 为 $C$ 的\Term{maximal-ideal}{极大理想}，对每个 $q>0$ 置
\[C_q=C/\mathfrak n^q,\qquad J_q=(J+\mathfrak n^q)/\mathfrak n^q.\]
于是 $C_q/J_q$ 识别为 $C/J$ 的一个商代数；复合\Term{homomorphism}{同态} $u_q:B\to C/J\to C_q/J_q$ 在目标赋\Term{discrete}{离散}拓扑时连续，且 $J_q$ 为 $C_q$ 的\Term{nilpotent}{幂零}理想。逐步构造 $A$-\Term{homomorphism}{同态} $v_q:B\to C_q$，使 (a) $v_q$ 提升 $u_q$；(b) $v_q$ 提升 $v_{q-1}$。归纳可行性容易验证，因为
\[u_q:B\to C/(J+\mathfrak n^q),\qquad v_{q-1}:B\to C/\mathfrak n^{q-1}\]
\SourcePage{71}
诱导同一个\Term{homomorphism}{同态}
\[B\to C/((J+\mathfrak n^q)+\mathfrak n^{q-1})=C/(J+\mathfrak n^{q-1})=C_{q-1}/J_{q-1},\]
即 $u_{q-1}$，所以它们定义\Term{homomorphism}{同态}
\[B\to C/J'_q,\qquad J'_q=(J+\mathfrak n^q)\cap\mathfrak n^{q-1}\supset\mathfrak n^q,\]
二者均由它\Term{reduced}{约化}得到。因此只需把 $B\to C/J'_q$ 提升到 $C_q$；这里 $J'_q/\mathfrak n^q$ 包含于 $J_q$，故\Term{nilpotent}{幂零}，由 (iii) 即可。这样，各 $v_q$ 定义从 $B$ 到各 $C_q$ 的\Term{projective-limit}{逆极限} $C$ 的\Term{homomorphism}{同态}。由于 $J$ 闭，它是各 $J_q$ 的\Term{projective-limit}{逆极限}，故 $v$ 提升 $u$。
(iv)$\Rightarrow$(i)。立即可见，若 (iv) 成立，则对任意\Term{finite}{有限}局部 $A$-代数 $A'$，$B\otimes_A A'$ 的各局部分量在 $A'$ 上也满足 (iv)。取在 $A$ 上自由且使 $B'$ 相对于 $A'$ 的各\Term{residue-field}{剩余域}\Term{group-extension}{扩张}\Term{trivial-bundle}{平凡}的 $A'$，由 1.4 (ii)，归结到 $B/A$ 的\Term{residue-field}{剩余域}\Term{group-extension}{扩张}\Term{trivial-bundle}{平凡}的情形。实际上我们将证明稍精细的结果如下。
\end{proof}
\begin{SGAresult}{推论}{2.2}
在 2.1 的条件下，再假定 $B/A$ 的\Term{residue-field}{剩余域}\Term{group-extension}{扩张}\Term{trivial-bundle}{平凡}。则 2.1 中的等价条件还等价于以下两个条件（在 (v) 中假定 $A,B$ 完备）：
\begin{enumerate}
\item[(iv ter)] 与 (iv) 相同，但在 $A$ 上\Term{finite}{有限}的 Artin \Term{local-ring}{局部环} $C$ 限为\Term{residue-field}{剩余域}\Term{group-extension}{扩张}\Term{trivial-bundle}{平凡}者（并且若愿意，还可要求 $J^2=0$）。
\item[(v)] 存在局部 $A$-\Term{homomorphism}{同态}
\[u:B\to B_1=A[[t_1,\ldots,t_n]],\qquad n=\dim\mathfrak n/(\mathfrak n^2+\mathfrak mB),\]
诱导\Term{isomorphism}{同构}
\[\mathfrak n/(\mathfrak n^2+\mathfrak mB)\xrightarrow{\sim}\mathfrak n_1/(\mathfrak n_1^2+\mathfrak mB_1),\]
其中 $\mathfrak n,\mathfrak n_1$ 分别为 $B,B_1$ 的\Term{maximal-ideal}{极大理想}，$\mathfrak m$ 为 $A$ 的\Term{maximal-ideal}{极大理想}。
\end{enumerate}
\end{SGAresult}
\begin{proof}
略去平方零理想这一附加小把戏，(iv bis) 显然蕴涵 (iv ter)，故只需证明 (iv ter)$\Rightarrow$(v)$\Rightarrow$(i)。
(iv ter)$\Rightarrow$(v)。取 $\mathfrak n/(\mathfrak n^2+\mathfrak mB)$ 的一组基 $a_1,\ldots,a_n$，便定义局部 $A$-代数\Term{homomorphism}{同态}
\[B\to B_1/(\mathfrak n_1^2+\mathfrak mB_1)=k[t_1,\ldots,t_n]/(t_1,\ldots,t_n)^2.\]
由 (iv ter)，逐步将其提升到 $B_1/\mathfrak n_1^2,B_1/\mathfrak n_1^3,\ldots$，取\Term{projective-limit}{逆极限}即得所需\Term{homomorphism}{同态} $B\to B_1$。
\SourcePage{72}
(v)$\Rightarrow$(i)。在\Term{commutative-diagram}{交换图}
\[\begin{tikzcd}[column sep=small]
\mathfrak m/\mathfrak m^2\arrow[r]\arrow[d]&\mathfrak n/\mathfrak n^2\arrow[r]\arrow[d]&\mathfrak n/(\mathfrak n^2+\mathfrak mB)\arrow[r]\arrow[d]&0\\
\mathfrak m/\mathfrak m^2\arrow[r]&\mathfrak n_1/\mathfrak n_1^2\arrow[r]&\mathfrak n_1/(\mathfrak n_1^2+\mathfrak mB_1)\arrow[r]&0
\end{tikzcd}\]
中，两行\Term{exact}{正合}，两端竖箭头\Term{surjectivity}{满射}，故中间竖箭头也\Term{surjectivity}{满射}；因 $B$ 完备，$B\to B_1$ \Term{surjectivity}{满射}。取 $x_i\in B$（$1\leq i\leq n$）提升 $t_i$，便定义 $A$-代数\Term{homomorphism}{同态} $B_1\to B$，它与 $u$ 一样因同一理由而\Term{surjectivity}{满射}，并由构造知其与 $u$ 的复合为恒等。因此 $B_1\to B$ 也\Term{injectivity}{单射}，从而是\Term{isomorphism}{同构}。于是得到：
\end{proof}
\begin{SGAresult}{推论}{2.3}
在 2.2 (v) 的条件下，$u$ 必为\Term{isomorphism}{同构}。
\end{SGAresult}
这就完成了 $B$ 在 $A$ 上\Term{formally-smooth}{形式光滑}的证明。同时也重新得到了 1.5（不过这实在算不上什么功劳）。
\SGAsection{3}{到光滑 S-概形的态射的局部无穷小延拓}
\begin{SGAresult}{定理}{3.1}
设 $f:X\to Y$ 为\Term{locally-finite-type}{局部有限型}\Term{morphism}{态射}。下列条件等价：
\begin{enumerate}
\item[(i)] $f$ \Term{smooth}{光滑}。
\item[(ii)] 对任意 $Y$ 上的\Term{prescheme}{预概形} $Y'$、与 $Y'$ 具有相同底空间的闭子\Term{prescheme}{预概形} $Y'_0$、任意 $Y$-\Term{morphism}{态射} $g_0:Y'_0\to X$ 及任意 $z\in Y'_0$，存在 $z$ 在 $Y'$ 中的\Term{open-neighborhood}{开邻域} $U$，使 $g_0|_{Y'_0\cap U}$ 可延拓为 $Y$-\Term{morphism}{态射} $g:U\to X$。
\item[(ii bis)] 对 (ii) 中的 $Y',Y'_0,z$，置 $X'=X\times_Y Y'$，$X'_0=X\times_Y Y'_0$；$X'_0$ 在 $Y'_0$ 上的任一\Term{section}{截面}，均可延拓为 $X'$ 在 $z$ 的某\Term{open-neighborhood}{开邻域} $U$ 上的\Term{section}{截面}。
\item[(iii)] 对任意 $Y$-\Term{scheme}{概形} $Y'$，若其为某个 $O_y$（$y\in Y$）上的\Term{finite}{有限} Artin \Term{local-ring}{局部环}的谱，则对 $Y'$ 的任意非空闭子\Term{prescheme}{预概形} $Y'_0$ 及任意 $Y$-\Term{morphism}{态射} $g_0:Y'_0\to X$，存在延拓 $g_0$ 的 $Y$-\Term{morphism}{态射} $g:Y'\to X$。
\item[(iii bis)] 对 (iii) 中的 $Y',Y'_0$，置 $X'=X\times_Y Y'$，$X'_0=X\times_Y Y'_0$；$X'_0$ 在 $Y'_0$ 上的任一\Term{section}{截面}，均可延拓为 $X'$ 在 $Y'$ 上的\Term{section}{截面}。
\end{enumerate}
\end{SGAresult}
\begin{proof}
(ii) 与 (ii bis)、(iii) 与 (iii bis) 的等价性以及 (ii)$\Rightarrow$(iii) 均显然。只需证明 (i)$\Rightarrow$(ii) 和 (iii)$\Rightarrow$(i)。
(i)$\Rightarrow$(ii)。置 $x=g_0(z)$。把 $X$ 换成 $x$ 的适当\Term{open-neighborhood}{开邻域}，并把 $Y'$ 换成该开集在 $g_0$ 下的原像上的诱导\Term{prescheme}{预概形}，便可假定 $X$ 在 $Y[t_1,\ldots,t_n]$ 上\Term{etale}{平展}。
\SourcePage{73}
复合 $Y$-\Term{morphism}{态射} $Y'_0\to X\to Y[t_1,\ldots,t_n]$ 由 $O_{Y'_0}$ 的 $n$ 个\Term{section}{截面}定义，这些\Term{section}{截面}可在 $z$ 的邻域内延拓为 $O_{Y'}$ 的\Term{section}{截面}，所以可假定上述\Term{morphism}{态射}已延拓为 $Y$-\Term{morphism}{态射} $Y'\to Y'_0$。由第 I 讲 5.6，存在唯一的 $Y$-\Term{morphism}{态射} $g:Y'\to X$ 提升这个\Term{morphism}{态射}并同时延拓 $g_0$。
(iii)$\Rightarrow$(i)。由于 $f$ \Term{smooth}{光滑}的点集开，只需证明它包含所有在其\Term{fiber}{纤维}中为\Term{closed-point}{闭点}的 $x\in X$。置 $y=f(x)$，则 $O_x$ 是某个\Term{finite-type}{有限型} $O_y$-代数的\Term{localization}{局部化}，\Term{residue-field}{剩余域}\Term{group-extension}{扩张}\Term{finite}{有限}。另一方面，(iii) 蕴涵：对任意\Term{finite}{有限} Artin 局部 $O_y$-代数 $A$ 及其\Term{jacobson-radical}{根基}中理想 $J$，每个\Term{homomorphism}{同态} $O_x\to A/J$ 均可提升为 $O_x\to A$。这里用到：从\Term{local-ring}{局部环} $B$ 的谱 $\operatorname{Spec}(B)$ 到 $X$ 的\Term{morphism}{态射}，与某个 $O_x$（$x\in X$）到 $B$ 的\Term{local-homomorphism}{局部同态}一一对应。由 2.1，$O_x$ 在 $O_y$ 上\Term{formally-smooth}{形式光滑}；由 1.9，它在 $O_y$ 上\Term{smooth}{光滑}。
\end{proof}
\begin{SGAresult}{推论}{3.2}
设 $f:X\to Y$ 如 3.1。下列条件等价：(i) $f$ \Term{etale}{平展}；(ii) 3.1 的条件 (ii) 成立，并且 $g_0$ 到 $U$ 的延拓 $g$ 唯一；(iii) 3.1 的条件 (iii) 成立，并且 $g$ 唯一。
\end{SGAresult}
只需注意，上述 (i)$\Rightarrow$(ii) 的证明中，只有当 $n=0$ 时才可能唯一（在 $z$ 的邻域内 $Y'_0$ 不等于 $Y'$ 时）；而这个条件已知也是充分的。
\begin{SGAresult}{推论}{3.3}
设 $X$ 为完备\Term{local-ring}{局部环} $A$ 上\Term{locally-finite-type}{局部有限型}\Term{prescheme}{预概形}，$y$ 为 $Y=\operatorname{Spec}(A)$ 的\Term{closed-point}{闭点}，$x$ 为 $f^{-1}(y)$ 中在 $k(y)$ 上有理的点。若 $X$ 在 $x$ 处相对于 $A$ \Term{smooth}{光滑}，则存在 $X$ 在 $Y$ 上“经过 $x$”的\Term{section}{截面} $s$，即 $s(y)=x$。
\end{SGAresult}
特别地，若 $X$ 在 $A$ 上\Term{smooth}{光滑}，则自然映射
\[\Gamma(X/Y)\longrightarrow\Gamma(X\otimes_A k/k)\]
从 $X$ 在 $Y$ 上的\Term{section}{截面}集到\Term{fiber}{纤维} $f^{-1}(y)=X\otimes_A k$ 的 $k$-有理点集是\Term{surjectivity}{满射}。这个事实在 $A$ 为\Term{discrete-valuation-ring}{离散赋值环}且 $X$ 在 $A$ 上\Term{proper}{固有}（实际上是\Term{projective}{射影}）的情形尤其熟知并得到应用；此时 $X$ 在 $Y$ 上的\Term{section}{截面}（即“$X$ 的 $A$-值点”）也可识别为有理\Term{section}{截面}，即 $X\otimes_A K=X_K$（它是 $K$ 上\Term{proper}{固有}且单纯的\Term{scheme}{概形}）的 $K$-值点，亦即 $X$ 的 $K$-有理点，其中 $K$ 是 $A$ 的\Term{field-of-fractions}{分式域}。
\SourcePage{74}
\SGAsection{4}{光滑 S-概形的局部无穷小延拓}
\begin{SGAresult}{定理}{4.1}
设 $Y$ 为\Term{locally-noetherian}{局部诺特}\Term{prescheme}{预概形}，$Y_0$ 为与它具有相同底空间的闭子\Term{prescheme}{预概形}，$X_0$ 为\Term{smooth}{光滑} $Y_0$-\Term{prescheme}{预概形}，$x\in X_0$。则存在 $x$ 的\Term{open-neighborhood}{开邻域} $U_0$、在 $Y$ 上\Term{smooth}{光滑}的\Term{prescheme}{预概形} $X$，以及 $Y_0$-\Term{isomorphism}{同构}
\[h:U_0\xrightarrow{\sim}X\times_Y Y_0.\]
此外，若 $(U'_0,X',h')$ 是问题的另一解，则“它在 $x$ 的邻域内与第一解\Term{isomorphism}{同构}”。
\end{SGAresult}
最后一句的确切含义留给读者。注意，给定 $U_0$ 后，问题的一个解等价于在 $U_0$ 上给出一个 $f_0^{-1}(O_Y)$-代数层 $\mathcal B$（$f_0$ 为结构\Term{morphism}{态射} $U_0\to Y_0$ 的底连续映射），以及环\Term{homomorphism}{同态} $\mathcal B\to O_{U_0}$，它与 $f^{-1}(O_Y)\to f^{-1}(O_{Y_0})$ 相容，并满足：(a) 该\Term{homomorphism}{同态}诱导\Term{isomorphism}{同构}
\[\mathcal B\otimes_{f^{-1}(O_Y)}f^{-1}(O_{Y_0})\xrightarrow{\sim}O_{Y_0};\]
(b) 在 $U_0$ 上赋予 $\mathcal B$ 后，它成为\Term{smooth}{光滑} $Y$-\Term{prescheme}{预概形}。这样，局部唯一性的确切意义就尤其清楚了。
\begin{proof}
可先假定 $X_0$ 在某个 $Y_0[t_1,\ldots,t_n]=Y'_0$ 上\Term{etale}{平展}。后者可视为 $Y'=Y[t_1,\ldots,t_n]$ 的闭子\Term{prescheme}{预概形}，且二者底空间相同。由第 I 讲 8.3，存在 $Y'$ 上\Term{etale}{平展}的 $X$ 及 $Y'_0$-\Term{isomorphism}{同构} $X\times_{Y'}Y'_0\xrightarrow{\sim}X'$。这就证明了存在性。唯一性由\Term{smooth}{光滑}\Term{morphism}{态射}的性质 3.1 (ii) 及下述引理得到。
\end{proof}
\begin{SGAresult}{引理}{4.2}
设 $Y$ 为\Term{prescheme}{预概形}，$Y_0$ 为由局部\Term{nilpotent}{幂零}\Term{ideal-sheaf}{理想层} $\mathcal J$ 定义的闭子\Term{prescheme}{预概形}；$X,X'$ 为 $Y$-\Term{prescheme}{预概形}，$u:X\to X'$ 为 $Y$-\Term{morphism}{态射}。假定 $X$ 在 $Y$ 上\Term{flat}{平坦}。则 $u$ 为\Term{isomorphism}{同构}当且仅当 $u_0:X\times_Y Y_0\to X'\times_Y Y_0$ 为\Term{isomorphism}{同构}。
\end{SGAresult}
化到\Term{affine}{仿射}情形，考察伴随\Term{grading}{分次}环，即容易证明。还要注意，把“\Term{isomorphism}{同构}”换成“\Term{closed-immersion}{闭浸入}”所得的对应命题也成立，而且不需要\Term{flat}{平坦}性假设。
\textbf{注记 4.3。}必须注意：4.1 所得的局部延拓 $X$ 不是典范的；换言之，两解之间的局部\Term{isomorphism}{同构}不唯一，即 $X$ 一般存在非\Term{trivial-bundle}{平凡}的 $Y$-\Term{automorphism}{自同构}，它们在闭子\Term{prescheme}{预概形} $X_0=X\times_Y Y_0$ 上诱导恒等。正因此，对于\Term{smooth}{光滑}\Term{prescheme}{预概形}的整体无穷小延拓，应当预期存在\Term{cohomology}{上同调}性质的障碍；下文第 6 节将予以说明。
\SourcePage{75}
\SGAsection{5}{态射的整体无穷小延拓}
设 $T$ 为\Term{topological-space}{拓扑空间}，$\mathcal G$ 为 $X$ 上的\Term{group-sheaf}{群层}，$\mathcal P$ 为 $T$ 上的\Term{set-sheaf}{集合层}，且 $\mathcal G$ 作用于 $\mathcal P$（为确定起见，取右作用）。若由此作用导出的熟知的\Term{set-sheaf}{集合层}\Term{homomorphism}{同态}
\[\mathcal G\times\mathcal P\longrightarrow\mathcal P\times\mathcal P\]
为\Term{isomorphism}{同构}，则称 $\mathcal P$ 在 $\mathcal G$ 下\Term{formally-principal-homogeneous}{形式主齐性}。等价地，对每个 $x\in T$，$\mathcal P_x$ 为空或为普通群 $\mathcal G_x$ 下的\Term{principal-homogeneous-space}{主齐性空间}；又等价于：对 $T$ 的每个开集 $U$，$\mathcal P(U)$ 为空或为普通群 $\mathcal G(U)$ 下的\Term{principal-homogeneous-space}{主齐性空间}。若 $\mathcal P$ \Term{formally-principal-homogeneous}{形式主齐性}，而且各 $\mathcal P_x$ 非空（即均为 $\mathcal G_x$ 下非空的\Term{principal-homogeneous-space}{主齐性空间}），则称 $\mathcal P$ 为 $\mathcal G$ 下的\Term{principal-homogeneous-sheaf}{主齐性层}。\footnote{似乎更宜采用 J. Giraud 在其博士论文中引入的、更简短明了的名称：“$\mathcal G$ 下的\Term{torsor}{挠子}”。}回顾：$\mathcal G$ 下\Term{principal-homogeneous-sheaf}{主齐性层}的\Term{isomorphism}{同构}类集识别为\Term{cohomology}{上同调}集合 $H^1(T,\mathcal G)$；当 $\mathcal G$ 交换时，它也是 $T$ 以 $\mathcal G$ 为系数的通常\Term{cohomology}{上同调}群。因而每个\Term{principal-homogeneous-sheaf}{主齐性层} $\mathcal P$ 对应特征类 $c(\mathcal P)\in H^1(T,\mathcal G)$；此类\Term{trivial-bundle}{平凡}当且仅当 $\mathcal P$ \Term{trivial-bundle}{平凡}（即\Term{isomorphism}{同构}于赋予右平移作用的 $\mathcal G$），又当且仅当 $\mathcal P$ 有\Term{section}{截面}。
\begin{SGAresult}{命题}{5.1}
设 $S$ 为\Term{prescheme}{预概形}，$X,Y$ 为 $S$ 上的\Term{prescheme}{预概形}，$Y_0$ 为由 $Y$ 上平方零\Term{ideal-sheaf}{理想层} $\mathcal J$ 定义的闭子\Term{prescheme}{预概形}。给定 $S$-\Term{morphism}{态射} $g_0:Y_0\to X$，令 $\mathcal P(g_0)$ 为 $Y$ 上如下的层：它在开集 $U$ 上的\Term{section}{截面}是把 $g_0|_{U\cap Y_0}$ 延拓为 $S$-\Term{morphism}{态射}的那些 $g:U\to X$。则 $\mathcal P(g_0)$ 自然是交换\Term{group-sheaf}{群层}
\[\mathcal G=\mathcal H\!om_{O_{Y_0}}(g_0^*(\Omega^1_{X/S}),\mathcal J)\]
下的\Term{formally-principal-homogeneous}{形式主齐性}层。
\end{SGAresult}
置 $\mathcal P=\mathcal P(g_0)$。需对每个 $Y$ 中开集 $U$ 定义映射 $\mathcal P(U)\times\mathcal G(U)\to\mathcal P(U)$，使：(a) 固定 $g\in\mathcal P(U)$ 时，$s\mapsto gs$ 是从 $\mathcal G(U)$ 到 $\mathcal P(U)$ 的\Term{bijection}{双射}；(b) $\mathcal P(U)$ 成为带 $\mathcal G(U)$ 作用的集合；(c) 这些映射与开集 $V\subset U$ 的限制算子相容。(c) 的验证显然，故可为简化而假定 $U=Y$。(b) 的验证（也可说具有局部性）留给读者。因此只需对给定的 $g\in\mathcal P(Y)$，定义从 $\mathcal G(Y)$ 到 $\mathcal P(Y)$ 的自然\Term{bijection}{双射}。
\SourcePage{76}
也就是说，已给定 $S$-\Term{morphism}{态射} $g:Y\to X$，要求典范\Term{bijection}{双射}
\[\operatorname{Hom}_{O_{Y_0}}(g_0^*(\Omega^1_{X/S}),\mathcal J)\xrightarrow{\sim}\mathcal P(Y).\tag{*}\]
这里 $\mathcal P(Y)$ 是那些与 $g$ 在 $Y_0$ 上诱导同一\Term{morphism}{态射} $g_0:Y_0\to X$ 的 $S$-\Term{morphism}{态射} $g':Y\to X$ 的集合。给定这样的 $g'$，等价于给定 $S$-\Term{morphism}{态射} $h:Y\to X\times_S X$，使 $\operatorname{pr}_1h=g$，$hi=(g_0,g_0)$；其中 $\operatorname{pr}_1:X\times_S X\to X$ 是第一投影，$i:Y_0\to Y$ 是典范\Term{immersion}{浸入}，$(g_0,g_0):Y_0\to X\times_S X$ 是分量为 $g_0,g_0$ 的\Term{morphism}{态射} $\Delta_{X/S}g_0$：
\[\begin{tikzcd}
X\times_S X\arrow[d,"\operatorname{pr}_1"']&Y_0\arrow[l,"{h_0=(g_0,g_0)}"']\arrow[d,"i"]\\
X&Y\arrow[l,"g"]\arrow[ul,dashed,"h"']
\end{tikzcd}\]
由于 $h_0$ 通过对角\Term{immersion}{浸入} $\Delta_{X/S}$ 分解，而 $Y$ 是 $Y_0$ 的一阶\Term{infinitesimal-neighborhood}{无穷小邻域}（即 $\mathcal J^2=0$），所求的 $h$ 必唯一地通过对角线的一阶\Term{infinitesimal-neighborhood}{无穷小邻域}分解。经第一投影将后者视为 $X$-\Term{prescheme}{预概形}，它可识别为代数层 $O_X+\Omega^1_{X/S}$ 的谱 $X'$（第二项视为平方零理想），\Term{diagonal-morphism}{对角态射} $X\to X'$ 对应这个代数层的典范增广。置 $Y'=X'\times_X Y$，$Y'_0=Y'\times_Y Y_0=X'\times_X Y_0$。于是所求的 $h$ 与 $Y'$ 在 $Y$ 上延拓给定\Term{section}{截面} $u_0$（$Y'_0$ 在 $Y_0$ 上的\Term{section}{截面}）的那些\Term{section}{截面} $u$ 一一对应。此外，$Y'$ 识别为 $Y$ 上代数层
\[\mathcal A=g^*(O_X+\Omega^1_{X/S})=O_Y+g^*(\Omega^1_{X/S})\]
的谱，$Y'_0$ 则对应代数层 $\mathcal A_0=\mathcal A\otimes_{O_Y}O_{Y_0}=O_{Y_0}+g_0^*(\Omega^1_{X/S})$；\Term{section}{截面} $u_0$ 由 $\mathcal A_0\to O_{Y_0}$ 的典范增广定义。因此 $\mathcal P(Y)$ 识别为那些诱导典范增广 $\mathcal A_0\to O_{Y_0}$ 的代数\Term{homomorphism}{同态} $\mathcal A\to O_Y$ 的集合。而这些代数\Term{homomorphism}{同态}与模\Term{homomorphism}{同态} $\mathcal M\to\mathcal A$ 一一对应，其中为简便置 $\mathcal M=g^*(\Omega^1_{X/S})$；我们关心其中诱导零\Term{homomorphism}{同态} $\mathcal M_0\to O_{Y_0}$ 的那些（$\mathcal M_0=\mathcal M\otimes_{O_Y}O_{Y_0}$），即把 $\mathcal M$ 映入\Term{augmentation-ideal}{增广理想} $\mathcal J$ 的那些。因此得到集合
\[\operatorname{Hom}_{O_Y}(\mathcal M,\mathcal J)=\operatorname{Hom}_{O_{Y_0}}(\mathcal M_0,\mathcal J),\]
因为 $\mathcal J$ 被 $\mathcal J$ 消去。这就是所求的典范\Term{bijection}{双射} (*)。
结合 3.1 中的蕴涵 (i)$\Rightarrow$(iii)，得到：
\begin{SGAresult}{推论}{5.2}
若 $X$ 在 $S$ 上\Term{smooth}{光滑}（至少在 $g_0(Y_0)$ 的各点处如此），则 $\mathcal P$ 实际上是交换\Term{group-sheaf}{群层} $\mathcal G$ 下的\Term{principal-homogeneous-sheaf}{主齐性层}，而且此时可写为
\[\mathcal G=g_0^*(\mathfrak g_{X/S})\otimes_{O_{Y_0}}\mathcal J.\]
\end{SGAresult}
\SourcePage{77}
其中 $\mathfrak g_{X/S}$ 为 $X$ 上 $\Omega^1_{X/S}$ 的对偶层，即 $X$ 相对于 $S$ 的\Term{tangent-sheaf}{切层}（或\Term{derivation}{导子}层）。（最后一个公式来自此时 $\Omega^1_{X/S}$ 为\Term{finite-type}{有限型}自由层这一事实。）特别地，这个\Term{principal-homogeneous-sheaf}{主齐性层}对应 $H^1(Y_0,\mathcal G)$ 中的\Term{cohomology}{上同调}类；此类为零是存在延拓 $g_0$ 的 $S$-\Term{morphism}{态射} $g$ 的充要条件。若存在一个这样的延拓，则所有延拓组成 $H^0(Y_0,\mathcal G)$ 下的齐性空间。
\Term{formal-geometry}{形式几何}方法的应用中，通常情形如下：给定两个 $S$-\Term{prescheme}{预概形} $X,Y$ 及 $S$ 上\Term{coherent}{凝聚}\Term{ideal-sheaf}{理想层} $\mathcal I$，以 $S_n$ 表示由 $\mathcal I^{n+1}$ 定义的闭子\Term{prescheme}{预概形}，置
\[X_n=X\times_S S_n,\qquad Y_n=Y\times_S S_n.\]
假定给定 $S_n$-\Term{morphism}{态射} $g_n:Y_n\to X_n$（等价地，$S$-\Term{morphism}{态射} $Y_n\to X$，也等价于 $S_{n+1}$-\Term{morphism}{态射} $Y_n\to X_{n+1}$，因为这样的\Term{morphism}{态射}必诱导 $Y_n\to X_n$），希望将其延拓为 $S_{n+1}$-\Term{morphism}{态射}
\[g_{n+1}:Y_{n+1}\to X_{n+1}.\]
（若能无限继续下去，就得到 $Y,X$ 分别沿\Term{ideal-sheaf}{理想层} $\mathcal I O_Y,\mathcal I O_X$ \Term{completion}{完备化}所得形式\Term{prescheme}{预概形}间的\Term{morphism}{态射} $\widehat Y\to\widehat X$。）可在 5.1 中把 $(S,X,Y,Y_0,g_0)$ 换成 $(S_{n+1},X_{n+1},Y_{n+1},Y_n,g_n)$。此处 $\mathcal G$ 成为从 $g_n^*(\Omega^1_{X_{n+1}/S_{n+1}})$ 到 $\mathcal J=\mathcal I^{n+1}O_Y/\mathcal I^{n+2}O_Y$ 的模\Term{homomorphism}{同态}层。由于 $\mathcal I O_Y$ 消去 $\mathcal J$，可将前一层换成其在 $Y_0$ 上诱导的层 $h_0^*(\Omega^1_{X/S})$，其中 $h_0$ 为复合 $Y_0\to Y_n\to X_{n+1}$，也即 $Y_0\to X_0\to X_{n+1}$，$g_0:Y_0\to X_0$ 由 $g_n$ 诱导。因 $\Omega^1_{X_{n+1}/S_{n+1}}$ 在 $X_0=X_{n+1}\times_{S_{n+1}}S_0$ 上的拉回为 $\Omega^1_{X_0/S_0}$，还有
\[\mathcal G=\mathcal H\!om_{O_{Y_0}}(g_0^*(\Omega^1_{X_0/S_0}),\mathcal I^{n+1}O_Y/\mathcal I^{n+2}O_Y).\]
于是得到：
\begin{SGAresult}{推论}{5.3}
设 $S,X,Y,\mathcal I,g_n$ 如上，$\mathcal P(g_n)$ 为 $Y$ 上的层，其在开集 $U$ 上的\Term{section}{截面}是把 $g_n$ 延拓为 $S_{n+1}$-\Term{morphism}{态射} $Y_{n+1}\to X_{n+1}$ 的那些 $g_{n+1}$。则 $\mathcal P(g_n)$ 是\Term{group-sheaf}{群层}
\[\mathcal G=\mathcal H\!om_{O_{Y_0}}(g_0^*(\Omega^1_{X_0/S_0}),\operatorname{gr}^{n+1}_{\mathcal I O_Y}(O_Y))\]
下的\Term{formally-principal-homogeneous}{形式主齐性}层。特别地：
\end{SGAresult}
\SourcePage{78}
\begin{SGAresult}{推论}{5.4}
若再假定 $X$ 在 $S$ 上\Term{smooth}{光滑}（至少在 $g_0(Y_0)$ 的各点如此），则 $\mathcal P(g_n)$ 实际上为\Term{principal-homogeneous-sheaf}{主齐性层}。特别地，它定义 $H^1(Y_0,\mathcal G)$ 中的\Term{obstruction-class}{障碍类}，其为零是存在 $g_n$ 的整体延拓 $g_{n+1}$ 的充要条件。若存在一个这样的延拓，则所有整体延拓组成 $H^0(Y_0,\mathcal G)$ 下的\Term{principal-homogeneous-space}{主齐性空间}。最后，此时\Term{group-sheaf}{群层}亦可写成
\[\mathcal G=g_0^*(\mathfrak g_{X_0/S_0})\otimes_{O_{Y_0}}\operatorname{gr}^{n+1}_{\mathcal I O_Y}(O_Y).\]
\end{SGAresult}
逐步进行便知：若所有 $H^1(Y_0,\mathcal G_n)$ 均为零，其中 $\mathcal G_n=g_0^*(\mathfrak g_{X_0/S_0})\otimes\operatorname{gr}^n_{\mathcal I O_Y}(O_Y)$，则从任意 $g_k$ 出发，均可依次延拓为 $g_{k+1},\ldots$。特别地，若 $\mathcal I$ \Term{nilpotent}{幂零}，就能把 $g_k$ 延拓到 $Y$。当 $Y_0$ \Term{affine}{仿射}时，这些 $H^1$ 特别都为零。因此：
\begin{SGAresult}{推论}{5.5}
在定理 3.1 中，把 (ii)（或 (ii bis)）里的 $Y'$ 假定为\Term{affine}{仿射}，并要求 $g_0$ 在整个 $Y'$ 上存在整体延拓 $g$，便得到一个与其他条件等价的充要条件。
\end{SGAresult}
注意 3.1 的证明不能直接给出这个结果。
一个重要情形是 $Y$ 在 $S$ 上\Term{flat}{平坦}。此时
\[\operatorname{gr}^n(O_Y)=\operatorname{gr}^n(O_S)\otimes_{O_{S_0}}O_{Y_0}.\]
若另有 $\operatorname{gr}^n(O_S)$ 在 $S$ 上\Term{locally-free}{局部自由}，则
\[\mathcal G_n=\mathcal H\!om_{O_{Y_0}}(g_0^*(\Omega^1_{X_0/S_0}),O_{Y_0})\otimes_{O_{S_0}}\operatorname{gr}^n(O_S);\]
若 $\Omega^1_{X_0/S_0}$ 也\Term{locally-free}{局部自由}（例如 $X$ 在 $S$ 上\Term{smooth}{光滑}），又可写为
\[\mathcal G_n=g_0^*(\mathfrak g_{X_0/S_0})\otimes_{O_{S_0}}\operatorname{gr}^n(O_S).\]
例如，$S$ 为\Term{affine}{仿射}，\Term{affine}{仿射}环为 $A$，$\mathcal I$ 由 $A$ 的理想 $I$ 定义时，对每个 $i$，
\[H^i(Y_0,\mathcal G_n)=H^i(Y_0,\mathcal G_0)\otimes_A\operatorname{gr}^n_I(A).\]
事实上问题在 $S_0$ 上局部，可归结到与\Term{free-module}{自由模}作\Term{tensor-product}{张量积}的情形。因此此时 $H^1(Y_0,\mathcal G_0)=0$ 蕴涵依次延拓 $g_n$ 的所有障碍均为零。故：
\begin{SGAresult}{推论}{5.6}
设 $(S,X,Y,\mathcal I,g_n)$ 如上，再假定 $X$ 在 $S$ 上\Term{smooth}{光滑}，$Y$ 在 $S$ 上\Term{flat}{平坦}，$S$ \Term{affine}{仿射}，且 $\operatorname{gr}^n(O_S)=\mathcal I^n/\mathcal I^{n+1}$ \Term{locally-free}{局部自由}。令 $\mathcal G_0=g_0^*(\mathfrak g_{X_0/S_0})$，则构造 $g_{n+1}$ 的障碍属于
\[H^1(Y_0,\mathcal G_0)\otimes_A\operatorname{gr}^{n+1}_I(A),\]
其中 $A$ 为 $S$ 的环，$I$ 为定义 $\mathcal I$ 的理想。若 $H^1(Y_0,\mathcal G_0)=0$，则 $g_n$ 可延拓为 $\widehat S$-\Term{morphism}{态射} $\widehat g:\widehat Y\to\widehat X$。
\end{SGAresult}
\SourcePage{79}
当然，若一开始以 $\widehat S$ 上 $\widehat{\mathcal I}$ 进形式\Term{prescheme}{预概形} $\mathfrak X,\mathfrak Y$ 代替普通 $S$-\Term{prescheme}{预概形} $X,Y$，结论也原样成立。例如可用它证明某些完备\Term{local-ring}{局部环}上的\Term{proper}{固有}形式\Term{scheme}{概形}实际上是代数的。确实，仿照引理 4.2 可得：
\begin{SGAresult}{推论}{5.7}
在 5.6 的条件下，若 $g_0$ 为\Term{isomorphism}{同构}，则 $\widehat g$ 亦然。
\end{SGAresult}
（注意：对于\Term{closed-immersion}{闭浸入}也有相同结论。）于是得到：
\begin{SGAresult}{命题}{5.8}
设 $A$ 为完备\Term{local-ring}{局部环}，\Term{maximal-ideal}{极大理想}为 $\mathfrak m$，\Term{residue-field}{剩余域}为 $k$；$\mathfrak X,\mathfrak Y$ 为 $A$ 上两个 $\mathfrak m$ 进形式\Term{prescheme}{预概形}，均在 $A$ 上\Term{flat}{平坦}（即对每个 $n$，$X_n,Y_n$ 均在 $A_n=A/\mathfrak m^{n+1}$ 上\Term{flat}{平坦}）。假定 $X_0=\mathfrak X\otimes_A k$ 在 $k$ 上\Term{smooth}{光滑}且 $H^1(X_0,\mathfrak g_{X_0/k})=0$。则每个从 $Y_0$ 到 $X_0$ 的 $k$-\Term{isomorphism}{同构}均可延拓为从 $\mathfrak Y$ 到 $\mathfrak X$ 的 $A$-\Term{isomorphism}{同构}；若另有 $H^0(X_0,\mathfrak g_{X_0/k})=0$，延拓唯一。
\end{SGAresult}
特别地，这就得到：在 $H^1(X_0,\mathfrak g_{X_0/k})=0$ 时，给定\Term{prescheme}{预概形} $X_0$ 的\Term{smooth}{光滑} $A$-形式\Term{prescheme}{预概形}提升具有唯一性。此外，若 $\mathfrak X,\mathfrak Y$ 来自 $A$ 上的普通\Term{proper}{固有}\Term{scheme}{概形} $X,Y$，则\Term{formal-geometry}{形式几何}中的\Term{existence-theorem}{层存在定理}（参见 Bourbaki 讨论班第 182 讲\footnote{证明参见 EGA III 5.4.1。}）说明，$A$-\Term{isomorphism}{同构} $Y\xrightarrow{\sim}X$ 与其形式\Term{completion}{完备化}的 $A$-\Term{isomorphism}{同构}一一对应。所以：
\begin{SGAresult}{推论}{5.9}
把 5.8 中的 $\mathfrak X,\mathfrak Y$ 换成 $A$ 上的普通\Term{proper}{固有}\Term{scheme}{概形} $X,Y$，命题仍成立。
\end{SGAresult}
最后，若 $\mathfrak X$ 为 $A$ 上\Term{proper}{固有}形式\Term{scheme}{概形}，而 $\mathfrak Y=\widehat Y$，$Y$ 是 $A$ 上普通\Term{proper}{固有}\Term{scheme}{概形}，那么 5.8 给出 $\mathfrak X$ \Term{isomorphism}{同构}于 $\widehat Y$ 的充分条件，从而给出 $\mathfrak X$ 实际上“代数”的充分条件；即它\Term{isomorphism}{同构}于某个 $\widehat X$，其中 $X$ 为 $A$ 上普通\Term{proper}{固有}\Term{scheme}{概形}，此时 $X$ 由它典范确定。特别当 $X_0=\mathbf P^r_k$（更一般地，$X_0$ 为 Severi--Brauer \Term{scheme}{概形}，即在 $k$ 的\Term{algebraically-closed}{代数闭}包上\Term{isomorphism}{同构}于标准\Term{projective}{射影}空间）时即是如此：$A$ 上任何\Term{proper}{固有}\Term{flat}{平坦}、\Term{fiber}{纤维}为 $\mathbf P^r_k$ 的形式\Term{scheme}{概形}均可\Term{algebraizable}{代数化}，更确切地说，\Term{isomorphism}{同构}于 $\mathbf P^r_A$ 的 $\mathfrak m$ 进形式\Term{completion}{完备化}。特别地（借助“存在定理”），$A$ 上\Term{fiber}{纤维}为 $\mathbf P^r_k$ 的普通\Term{proper}{固有}\Term{scheme}{概形}均\Term{isomorphism}{同构}于 $\mathbf P^r_A$（$A$ 为完备\Term{local-ring}{局部环}）。借助\Term{descent}{下降}理论可证明：若 $A$ 不完备，则作\Term{finite}{有限}\Term{etale}{平展}\Term{base-extension}{基扩张} $A\to A'$ 后，$X$ \Term{isomorphism}{同构}于 $\mathbf P^r$；在这种形式下，\Term{fiber}{纤维}为 Severi--Brauer \Term{scheme}{概形}时结论仍成立。
\SourcePage{80}
\SGAsection{6}{光滑 S-概形的整体无穷小延拓}
在定理 4.1 的条件下，我们要问：是否存在 $Y$ 上\Term{smooth}{光滑}\Term{prescheme}{预概形} $X$，使 $X\times_Y Y_0$ 与 $X_0$ 在 $Y_0$ 上\Term{isomorphism}{同构}？已知这样的\Term{scheme}{概形}“在 $X_0$ 上局部存在”。再次采用逐步构造的方法，将字母 $Y$ 换为 $S$，假定给定 $S$ 中由\Term{ideal-sheaf}{理想层} $\mathcal I$ 定义的闭子\Term{prescheme}{预概形} $S_0$（不再需要假定 $\mathcal I$ 局部\Term{nilpotent}{幂零}），引入由各 $\mathcal I^{n+1}$ 定义的闭子\Term{prescheme}{预概形} $S_n$，并假定给定在 $S_n$ 上\Term{smooth}{光滑}的子\Term{prescheme}{预概形} $X_n$。希望找到“\Term{reduced}{约化}为 $X_n$”的 $S_{n+1}$-\Term{prescheme}{预概形} $X_{n+1}$，即带有\Term{isomorphism}{同构}
\[X_{n+1}\times_{S_{n+1}}S_n\xleftarrow{\sim}X_n,\]
而且它在 $S_{n+1}$ 上\Term{smooth}{光滑}（由第 II 讲 2.1，等价于在 $S_{n+1}$ 上\Term{flat}{平坦}）。如第 4 节所说，这等价于给定 $f^{-1}(O_{S_{n+1}})$ 上代数层 $\mathcal B$（$f$ 为结构\Term{morphism}{态射} $X_n\to S_n$ 的底连续映射），带有与增广 $f^{-1}(O_{S_{n+1}})\to f^{-1}(O_{S_n})$ 相容的增广 $\mathcal B\to O_{X_n}$，并满足两个条件 (a)、(b)。这里不再重写它们，只指出这些条件在 $X_n$ 的底\Term{topological-space}{拓扑空间}上是局部的。由 4.1 知解局部存在，而且至少局部在\Term{isomorphism}{同构}意义下唯一，但\Term{isomorphism}{同构}不唯一。先把这一点说清楚。
\begin{SGAresult}{命题}{6.1}
设 $X_{n+1}$ 在 $S_{n+1}$ 上，\Term{reduced}{约化}为 $S_n$ 上的 $X_n$。则 $X_{n+1}$ 的那些在 $X_n$ 上诱导恒等的 $S_{n+1}$-\Term{automorphism}{自同构}所成的层（定义于 $X_n$，亦即 $X_0$ 的底\Term{topological-space}{拓扑空间}上），典范\Term{isomorphism}{同构}于\Term{group-sheaf}{群层}
\[\mathcal G=\mathfrak g_{X_0/S_0}\otimes_{O_{S_0}}\operatorname{gr}^{n+1}_{\mathcal I}(O_S).\]
\end{SGAresult}
\begin{proof}
由 5.4 和 4.2，这个层是 $\mathcal G$ 下的\Term{principal-homogeneous-sheaf}{主齐性层}。它带有一个特选\Term{section}{截面}，即 $X_{n+1}$ 的恒等\Term{automorphism}{自同构}，故作为\Term{set-sheaf}{集合层}可识别为 $\mathcal G$。还需验证这个识别与群结构相容。这很容易，事实上是更一般结果的特例：5.1、5.3 中\Term{principal-bundle}{主丛}结构与\Term{morphism}{态射}复合相容（这里不陈述，但这应收入那部巨著）。
\end{proof}
特别地，$X_0$ 上 $X_{n+1}$ 的\Term{automorphism}{自同构}芽层（带有上述结构）交换。因此，若 $X'_{n+1}$ 为另一解，并在 $X_0$ 的开集 $U$ 上\Term{isomorphism}{同构}于 $X_{n+1}$，则由\Term{isomorphism}{同构} $X_{n+1}|_U\xrightarrow{\sim}X'_{n+1}|_U$ 搬移结构所得的
\[\operatorname{Aut}(X_{n+1})|_U\longrightarrow\operatorname{Aut}(X'_{n+1})|_U\]
不依赖于该\Term{isomorphism}{同构}的选择。
\SourcePage{81}
（事实上，若由 6.1 将两个\Term{automorphism}{自同构}层都识别为 $\mathcal G$，上述\Term{isomorphism}{同构}正是 $\mathcal G$ 的恒等\Term{isomorphism}{同构}。）由 6.1 得：
\begin{SGAresult}{推论}{6.2}
设 $X_{n+1},X'_{n+1}$ 在 $S_{n+1}$ 上\Term{smooth}{光滑}并“\Term{reduced}{约化}为 $X_n$”。则从 $X_{n+1}$ 到 $X'_{n+1}$ 的那些在 $X_n$ 上诱导恒等的 $S_{n+1}$-\Term{isomorphism}{同构}组成的层（在 $X_0$ 的底空间上），自然是 $\mathcal G$ 下的\Term{principal-homogeneous-sheaf}{主齐性层}。
\end{SGAresult}
这确实只是说：$X_{n+1},X'_{n+1}$ 局部\Term{isomorphism}{同构}，而且前者的\Term{automorphism}{自同构}芽层为 $\mathcal G$。
现在注意，由 4.1，总能取 $X_n$ 的开\Term{covering}{覆盖} $(U_i)$（可令各 $U_i$ \Term{affine}{仿射}），并为每个 $i$ 取 $S_{n+1}$ 上\Term{smooth}{光滑}\Term{scheme}{概形} $X^i$，\Term{reduced}{约化}为 $U_i$。为简便假定 $X_n$ \Term{separated}{分离}，则 $U_{ij}=U_i\cap U_j$ 仍是 $X_n$ 中\Term{affine}{仿射}开集。由于这种开集以\Term{quasi-coherent}{拟凝聚}层 $\mathcal G$ 为系数的 $H^1$ 为零，由 6.2 知 $X^i|_{U_{ij}}$ 与 $X^j|_{U_{ij}}$ \Term{isomorphism}{同构}。取一个这样的\Term{isomorphism}{同构}
\[f_{ji}:X^i|_{U_{ij}}\xrightarrow{\sim}X^j|_{U_{ij}},\]
它在相差一个 $\mathcal G$ 在 $U_{ij}$ 上\Term{section}{截面}的意义下确定。对每组三重指标，置
\[f_{ji}^{(k)}=f_{ji}|_{U_{ijk}},\qquad U_{ijk}=U_i\cap U_j\cap U_k.\]
若有
\[f_{kj}^{(i)}f_{ji}^{(k)}=f_{ki}^{(j)},\tag{1}\]
则各 $X^i$ 可沿 $f_{ji}$ “\Term{gluing}{粘合}”，从而定义所求问题的解 $X=X_{n+1}$。更一般地，若可把各 $f_{ji}$ 修改为
\[f'_{ji}=f_{ji}g_{ji},\qquad g_{ji}\in\Gamma(U_{ij},\mathcal G),\tag{2}\]
使 $f'_{ji}$ 满足上述传递条件，也就有解。这个充分条件也必要：回顾任何这样的解 $X$ 在每个 $U_i$ 上都须\Term{isomorphism}{同构}于 $X^i$，故可取\Term{isomorphism}{同构} $f_i:X|_{U_i}\xrightarrow{\sim}X^i$，再定义
\[f'_{ji}=(f_j|_{U_{ij}})(f_i|_{U_{ij}})^{-1}:X^i|_{U_{ij}}\xrightarrow{\sim}X^j|_{U_{ij}},\]
它们满足\Term{gluing}{粘合}条件。现置
\[f_{ijk}=(f_{ki}^{(j)})^{-1}f_{kj}^{(i)}f_{ji}^{(k)}.\tag{3}\]
\SourcePage{82}
这是 $X^i|_{U_{ijk}}$ 的\Term{automorphism}{自同构}；借助 6.1，将它识别为 $\mathcal G$ 的一个\Term{section}{截面}。作一个留给读者的小型形式计算可见，它是开\Term{covering}{覆盖} $\mathcal U=(U_i)$ 以 $\mathcal G$ 为系数的一个 $2$-上\Term{cycle}{闭链} $f$。同一计算表明，在 (2) 下，$f'_{ij}$ 的\Term{gluing}{粘合}条件 (1) 等价于
\[f=dg,\tag{4}\]
其中把 $g=(g_{ij})$ 看作 $\mathcal U$ 以 $\mathcal G$ 为系数的 $1$-上链。因此，问题有解的充要条件是：上\Term{cycle}{闭链} (3) 所定义的 $H^2(\mathcal U,\mathcal G)$ 中的\Term{cohomology}{上同调}类为零。又因 $\mathcal U$ 是\Term{scheme}{概形} $X_0$ 的\Term{affine}{仿射}\Term{covering}{覆盖}，$H^2(\mathcal U,\mathcal G)$ 识别为 $H^2(X_0,\mathcal G)$。所得的\Term{cohomology}{上同调}类显然不依赖于\Term{affine}{仿射}\Term{covering}{覆盖}的选择，称为把 $X_n$ 延拓为 $S_{n+1}$ 上\Term{smooth}{光滑}\Term{scheme}{概形} $X_{n+1}$ 的\Term{obstruction-class}{障碍类}。
假定这个障碍为零。由上文略述的推理，任何解 $X=X_{n+1}$ 均\Term{isomorphism}{同构}于利用\Term{isomorphism}{同构} $f'_{ji}$ \Term{gluing}{粘合}所得的一个解；可假定 $f'_{ji}$ 形如 (2)，而\Term{gluing}{粘合}条件正是 (3)。于是所有可取的 $g$ 组成 $\mathcal U$ 以 $\mathcal G$ 为系数的 $1$-上\Term{cycle}{闭链}群 $Z^1(\mathcal U,\mathcal G)$ 下的\Term{principal-homogeneous-space}{主齐性空间}。此外，立刻可见：两个上链 $g,g'$（满足 $dg=dg'=f$）定义\Term{isomorphism}{同构}的解，当且仅当上\Term{cycle}{闭链} $g-g'$ 形如 $dh$，其中 $h=(h_i)\in C^0(\mathcal U,\mathcal G)$。因此：
\begin{SGAresult}{定理}{6.3}
设 $(S,\mathcal I,X_n)$ 如上，假定 $X_n$ \Term{separated}{分离}。\footnote{这一条件实际上不必要，也可以避开上面的上\Term{cycle}{闭链}计算。参见 J. Giraud 的书 \emph{Cohomologie Non Abélienne}（将于 1971 年由 Springer-Verlag 出版）。比较注记 6.4。}则可典范定义 $H^2(X_0,\mathcal G)$ 中的\Term{obstruction-class}{障碍类}（$\mathcal G$ 如 6.1），其为零是存在 $S_{n+1}$ 上\Term{smooth}{光滑}且\Term{reduced}{约化}为 $X_n$ 的\Term{scheme}{概形} $X_{n+1}$ 的充要条件。若障碍为零，则\Term{reduced}{约化}为 $X_n$ 的 $S_{n+1}$-\Term{prescheme}{预概形} $X_{n+1}$ 的\Term{isomorphism}{同构}类集（\Term{isomorphism}{同构}要求在 $X_n$ 上诱导恒等），自然为 $H^1(X_0,\mathcal G)$ 下的\Term{principal-homogeneous-space}{主齐性空间}。
\end{SGAresult}
\textbf{注记 6.4。}自 6.1 起的推理完全是形式的，放到局部\Term{category}{范畴}乃至一般\Term{fibered-category}{纤维范畴}的框架中会更为自然。在一般背景下，一个\Term{category}{范畴}的“整体”对象的存在\Term{obstruction-class}{障碍类}，可在该\Term{category}{范畴}局部有对象、任意两个对象“局部\Term{isomorphism}{同构}”且每个对象的\Term{automorphism}{自同构}群交换时如此得到。它包含如下特例：不必交换的\Term{group-sheaf}{群层}\Term{exact-sequence}{正合列}中的“第二连接\Term{homomorphism}{同态}”（例如 Grothendieck 在 Kansas 或 Tôhoku 中所研究的）。
\SourcePage{83}
因此，这里所作的直接上\Term{cycle}{闭链}计算，只应视为缺乏适当参考文献时的权宜办法。
\textbf{6.5。}注意，6.3 的 $H^1(X_0,\mathcal G)$ 下\Term{principal-homogeneous-space}{主齐性空间}一般没有特选元素。特别地，在 $S$ 上\Term{localization}{局部化}后，得到 $S_0$ 上结构群为 $R^1f_*(\mathcal G)$ 的\Term{principal-homogeneous-sheaf}{主齐性层}，它不必\Term{trivial-bundle}{平凡}，即定义 $H^1(S_0,R^1f_*(\mathcal G))$ 中不必为零的\Term{cohomology}{上同调}类。（这里假定 $d\in H^2(X_0,\mathcal G)$ 不为零，但“在 $S$ 上局部”为零，即定义 $R^2f_*(\mathcal G)$ 的零\Term{section}{截面}，亦即 $H^0(S_0,R^2f_*(\mathcal G))$ 中的零元素。）
\textbf{6.6。}关于本节引入的\Term{cohomology}{上同调}类的一般代数机制，以及它们与上一节的联系，目前几乎一无所知；即使在最简单的特例，例如 Artin 环上的阿贝尔\Term{scheme}{概形}情形，也无法作出确切说明。\footnote{现在已经知道，这个情形中的障碍总为零。[2003 年增补（MR）：参见 F. Oort, \emph{Finite group schemes, local moduli for abelian varieties and lifting problems}, Algebraic Geometry Oslo 1970, Wolters-Noordhoff, 1972, pp. 223--254。]}希望有人深入钻研这个看来十分有趣的问题。它特别与代数结构的“模理论”密切相关。
\begin{SGAresult}{推论}{6.7}
若 $H^2(X_0,\mathcal G)=0$，则 $X_{n+1}$ 存在；若再有 $H^1(X_0,\mathcal G)=0$，则它在\Term{isomorphism}{同构}意义下唯一。
\end{SGAresult}
逐步应用，并注意\Term{affine}{仿射}\Term{scheme}{概形}对\Term{quasi-coherent}{拟凝聚}层无高阶\Term{cohomology}{上同调}，特别得到：
\begin{SGAresult}{推论}{6.8}
在定理 4.1 的条件下，若 $X_0$ \Term{affine}{仿射}，则存在 $Y$ 上\Term{smooth}{光滑}且\Term{reduced}{约化}为 $X_0$ 的 $X$；这个 $X$ 在\Term{isomorphism}{同构}意义下唯一，但\Term{isomorphism}{同构}不唯一。
\end{SGAresult}
注意定理 4.1 的直接证明不能给出这一结果。
\begin{SGAresult}{推论}{6.9}
在 6.3 的条件下，假定 $S$ \Term{affine}{仿射}，环为 $A$，$\mathcal I$ 由 $A$ 的理想 $I$ 定义，且 $\operatorname{gr}^n_I(O_S)=\mathcal I^n/\mathcal I^{n+1}$ \Term{locally-free}{局部自由}。则
\[H^i(X_0,\mathcal G)=H^i(X_0,\mathcal G_0)\otimes_A\operatorname{gr}^{n+1}_I(A),\qquad\mathcal G_0=\mathfrak g_{X_0/S_0}.\]
因此延拓 $X_n$ 的\Term{obstruction-class}{障碍类}属于 $H^2(X_0,\mathcal G_0)\otimes_A\operatorname{gr}^{n+1}_I(A)$；若它为零，则解的\Term{isomorphism}{同构}类集是 $H^1(X_0,\mathcal G_0)\otimes_A\operatorname{gr}^{n+1}_I(A)$ 下的\Term{principal-homogeneous-space}{主齐性空间}。
\end{SGAresult}
\SourcePage{84}
特别地：
\begin{SGAresult}{推论}{6.10}
在 6.9 的条件下，假定 $H^2(X_0,\mathfrak g_{X_0/S_0})=0$。则在 $S$ 的 $\mathcal I$ 进形式\Term{completion}{完备化} $\widehat S$ 上，存在 $\widehat{\mathcal I}$ 进形式\Term{scheme}{概形} $\mathfrak X$，它“在 $S$ 上\Term{smooth}{光滑}”（即各 $X_p$ 在 $S_p$ 上\Term{smooth}{光滑}），并\Term{reduced}{约化}为 $X_n$，即带有\Term{isomorphism}{同构}
\[\mathfrak X\times_S S_n\xleftarrow{\sim}X_n.\]
若另有 $H^1(X_0,\mathfrak g_{X_0/S_0})=0$，则这样的 $\mathfrak X$ 在\Term{isomorphism}{同构}意义下唯一。
\end{SGAresult}
确实，可依次构造 $X_{n+1},X_{n+2},\ldots$，再取各 $X_i$ 的正向极限得到 $\mathfrak X$。唯一性已在上一节证明。
\SGAsection{7}{应用：在完备局部环 A 上构造光滑形式概形及普通概形}
上一节的结果有时可以证明：给定 $k$ 上\Term{smooth}{光滑}\Term{scheme}{概形} $X_0$，存在上述环上的 $\mathfrak m$ 进形式\Term{scheme}{概形}\Term{reduced}{约化}为 $X_0$。分两种情形讨论。
\textbf{a) $A$ 为“等特征”。}当 $k$ 的特征为 $0$ 时尤其如此。已知 $A$ 有代表子域，即存在子域 $k'$，使 $A\to k$ 诱导\Term{isomorphism}{同构} $k'\xrightarrow{\sim}k$。此时甚至存在\Term{reduced}{约化}为 $X_0$ 的普通\Term{smooth}{光滑} $A$-\Term{scheme}{概形}，即 $X=X_0\otimes_k A$；这里通过 $k'$ 所定义的\Term{homomorphism}{同态} $k\to k'\to A$，把 $A$ 视为 $k$-代数。但是这个构造不是“自然”的：容易看出，即使 $A=k[t]/(t^2)$ 为对偶数代数时，换一个提升\Term{homomorphism}{同态} $k\to A$（此时由 $k$ 到自身的绝对\Term{derivation}{导子}定义），所得的 $A$-\Term{scheme}{概形} $X'$ 一般也不\Term{isomorphism}{同构}于 $X$（若 $H^1(X_0,\mathfrak g_{X_0/k})\ne0$）。当 $k$ 特征为 $0$，或特征为 $p>0$ 且不\Term{perfect}{完美}时，研究这样得到哪些\Term{smooth}{光滑} $A$-\Term{scheme}{概形} $X$、两个\Term{homomorphism}{同态} $k\to A$ 在什么条件下定义\Term{isomorphism}{同构}的 $A$-\Term{scheme}{概形}，会很有趣。不过 $k'$ 的存在已足以保证提升 $X_0$ 的第一障碍
\[H^2(X_0,\mathfrak g_{X_0/k})\otimes_k\mathfrak m/\mathfrak m^2\]
中的类必为零。当然，把 $X_0$ 提升为 $A/\mathfrak m^2$ 上\Term{smooth}{光滑} $X_1$ 后，构造 $X_2$ 的新障碍一般不为零：它是 $H^1(X_0,\mathcal G_0)\otimes\mathfrak m/\mathfrak m^2$ 下某个\Term{principal-homogeneous-space}{主齐性空间}中可变元素的函数，取值于 $H^2(X_0,\mathcal G_0)\otimes\mathfrak m^2/\mathfrak m^3$；应当详细研究这一情况。\footnote{它很可能由 Kodaira--Spencer 括号运算描述（参见 Cartan 讨论班，1960/61，第 4 讲）。}
\SourcePage{85}
\textbf{b) $A$ 为混合特征。}此时我们一无所知，除非碰巧 $H^2(X_0,\mathfrak g_{X_0/k})=0$，这时可在 $A$ 上构造\Term{reduced}{约化}为 $k$ 的单纯 $\mathfrak m$ 进形式\Term{scheme}{概形}。即使 $A=\mathbf Z/p^2\mathbf Z$ 且 $X_0$ 是二维“阿贝尔”\Term{scheme}{概形}，也不知道能否将它提升为 $A$ 上\Term{smooth}{光滑}的 $X=X_1$。\footnote{现在已经证明可以，参见原书第 67 页注 7。}另一方面，也没有已经证明不来自 $A$ 上普通\Term{smooth}{光滑}\Term{scheme}{概形} $X$ 的 $X_0$ 的例子。（我觉得应当存在这样的例子，且 $X_0$ 可取为\Term{projective}{射影}曲面。）\footnote{后来 J.-P. Serre 构造了这样的例子（\emph{Proc. Nat. Acad. Sc. USA} 47 (1961), no. 1, pp. 108--109），至少在某些维数如此。D. Mumford 给出了代数曲面的例子（未发表）。}这里只指出：由 Cohen 定理，存在\Term{residue-field}{剩余域}为 $k$ 的 Cohen $p$-环 $B$，以及在\Term{residue-field}{剩余域}上诱导恒等\Term{isomorphism}{同构}的\Term{homomorphism}{同态} $B\to A$；所以取 $A$ 为 Cohen $p$-环便可得到“最强”的提升结果：若在这样的环上有普通或形式解，则在任何\Term{residue-field}{剩余域}为 $k$ 的完备\Term{local-ring}{局部环}上都有解。特别地，对于 Cohen $p$-环，$\mathfrak m/\mathfrak m^2$ 典范识别为 $k$，因此对于特征 $p>0$ 的域 $k$ 上任意\Term{smooth}{光滑}\Term{scheme}{概形} $X_0$，有一个 $H^2(X_0,\mathfrak g_{X_0/k})$ 中的\Term{cohomology}{上同调}类，作为将 $X_0$ 提升为 Cohen $p$-环上\Term{smooth}{光滑}\Term{scheme}{概形}的第一障碍；我们不知道它能否非零。\footnote{它确实可以非零，如原书第 69 页注 10 所述。}
即使能逐步构造\Term{reduced}{约化}为 $X_0$ 的各 $X_n$，一般也只得到 $A$ 上\Term{smooth}{光滑}且\Term{reduced}{约化}为 $X_0$ 的形式\Term{scheme}{概形} $\mathfrak X$。当 $X_0$ 在 $A$ 上\Term{proper}{固有}时，为得到 $A$ 上普通\Term{proper}{固有}且单纯、\Term{reduced}{约化}为 $X_0$ 的\Term{scheme}{概形}，还要问 $\mathfrak X$ 是否\Term{algebraizable}{可代数化}。唯一已知的判据（Bourbaki 讨论班中提及，载于《几何基础》第 III 章 4.7.1）如下：若 $\mathfrak X$ 在 $A$ 上\Term{proper}{固有}，$\mathcal L$ 为 $\mathfrak X$ 上\Term{invertible-sheaf}{可逆层}，且其在 $X_0$ 上诱导的 $\mathcal L_0$ \Term{ample}{丰沛}（即某个张量幂 $\mathcal L_0^{\otimes n}$，$n>0$，来自 $X_0$ 的\Term{projective}{射影}\Term{immersion}{浸入}），则存在 $A$ 上\Term{projective}{射影}\Term{scheme}{概形} $X$ 及 $X$ 上\Term{ample}{丰沛}\Term{invertible-sheaf}{可逆层}，使 $(\mathfrak X,\mathcal L)$ 由其作 $\mathfrak m$ 进\Term{completion}{完备化}得到。这就引出：给定 $X_0$ 上\Term{locally-free}{局部自由}层 $\mathcal E_0$（为上述目的将取为\Term{ample}{丰沛}\Term{invertible-sheaf}{可逆层}），将它延拓为 $\mathfrak X$ 上\Term{locally-free}{局部自由}层 $\mathcal E$。只需逐步在各 $X_n$ 上构造\Term{locally-free}{局部自由}层 $\mathcal E_n$。讨论与第 6 节完全类似（参见注记 6.4）；关键是 $\mathcal E_{n+1}$ 的那些在 $\mathcal E_n$ 上诱导恒等的\Term{automorphism}{自同构}所组成的层。立即证明它识别为
\[\mathcal G=\mathcal H\!om_{O_{X_0}}(\mathcal E_0,\mathcal E_0\otimes\operatorname{gr}^{n+1}_{\mathcal I}(O_X))=\mathcal H\!om_{O_{X_0}}(\mathcal E_0,\mathcal E_0)\otimes\operatorname{gr}^{n+1}_{\mathcal I}(O_X).\tag{*}\]
它仍是交换\Term{group-sheaf}{群层}。得到：
\SourcePage{86}
\begin{SGAresult}{命题}{7.1}
设 $S$ 为\Term{prescheme}{预概形}，带有\Term{quasi-coherent}{拟凝聚}\Term{ideal-sheaf}{理想层} $\mathcal I$；$X$ 为 $S$ 上\Term{prescheme}{预概形}，$S_n$ 为 $\mathcal I^{n+1}$ 定义的子\Term{prescheme}{预概形}，$X_n=X\times_S S_n$（对每个整数 $n$）。设 $\mathcal E_n$ 为 $X_n$ 上\Term{locally-free}{局部自由}层，希望将其延拓为 $X_{n+1}$ 上\Term{locally-free}{局部自由}层 $\mathcal E_{n+1}$。则 $\mathcal E_n$ 定义 $H^2(X_0,\mathcal G)$ 中的典范\Term{obstruction-class}{障碍类}，$\mathcal G$ 为上式给出的\Term{quasi-coherent}{拟凝聚}层；此类为零是存在延拓 $\mathcal E_n$ 的 $\mathcal E_{n+1}$ 的充要条件。若它为零，则解 $\mathcal E_{n+1}$ 的\Term{isomorphism}{同构}类集（\Term{isomorphism}{同构}在 $\mathcal E_n$ 上诱导恒等）是 $H^1(X_0,\mathcal G)$ 下的\Term{principal-homogeneous-space}{主齐性空间}。
\end{SGAresult}
这个命题给出通常的那些推论。这里只指出：若 $X$ 在 $S$ 上\Term{flat}{平坦}，则可写为
\[\mathcal G=\mathcal H\!om_{O_{X_0}}(\mathcal E_0,\mathcal E_0)\otimes_{O_{S_0}}\operatorname{gr}^{n+1}_{\mathcal I}(O_S).\]
因此，若 $S$ \Term{affine}{仿射}、环为 $A$，且各 $\mathcal I^n/\mathcal I^{n+1}$ \Term{locally-free}{局部自由}，则 $\mathcal E_{n+1}$ 存在的一个充分条件为
\[H^2(X_0,\mathcal G_0)=0,\qquad\mathcal G_0=\mathcal H\!om_{O_{X_0}}(\mathcal E_0,\mathcal E_0).\]
于是逐步可得到相继的延拓 $\mathcal E_m$（$m=n,n+1,\ldots$）。回到起初的情形，得到：
\begin{SGAresult}{命题}{7.2}
设 $A$ 为完备\Term{local-ring}{局部环}，$\mathfrak X$ 为 $A$ 上\Term{proper}{固有}\Term{flat}{平坦}形式\Term{scheme}{概形}，$X_0$ \Term{projective}{射影}且 $H^2(X_0,O_{X_0})=0$。则存在 $A$ 上\Term{projective}{射影}\Term{scheme}{概形} $X$，其 $\mathfrak m$ 进形式\Term{completion}{完备化}\Term{isomorphism}{同构}于 $\mathfrak X$。
\end{SGAresult}
结合 6.10，得到：
\begin{SGAresult}{定理}{7.3}
设 $A$ 为\Term{residue-field}{剩余域}为 $k$ 的完备\Term{local-ring}{局部环}，$X_0$ 为 $k$ 上\Term{projective}{射影}\Term{smooth}{光滑}\Term{scheme}{概形}，且
\[H^2(X_0,\mathfrak g_{X_0/k})=H^2(X_0,O_{X_0})=0.\]
则存在 $A$ 上\Term{smooth}{光滑}\Term{projective}{射影}\Term{scheme}{概形} $X$，\Term{reduced}{约化}为 $X_0$。
\end{SGAresult}
更一般地，若给定在 $A_n=A/\mathfrak m^{n+1}$ 上\Term{smooth}{光滑}且\Term{reduced}{约化}为 $X_0$ 的 $X_n$，则存在 $A$ 上\Term{smooth}{光滑}\Term{proper}{固有}的 $X$ 及\Term{isomorphism}{同构} $X\otimes_A A_n=X_n$。
\begin{SGAresult}{推论}{7.4}
$k$ 上每条\Term{smooth}{光滑}\Term{proper}{固有}曲线，均可由 $A$ 上\Term{smooth}{光滑}\Term{proper}{固有}曲线\Term{reduced}{约化}得到。
\end{SGAresult}
这一结果将与\Term{formal-geometry}{形式几何}中的\Term{existence-theorem}{层存在定理}一起，成为以超越方法研究 $X_0$ 的\Term{fundamental-group}{基本群}的主要工具。
